\subsection{函数}

定义9. 若对于每个x，在F下至多有一个对象与x对应（第一章，§5，第3节），则称图F为函数图。一个对应 \(f = (\mathbf{F}, \mathbf{A}, \mathbf{B})\) 被称为函数，如果其图像 F 是函数图像，并且其源 A 等于其定义域 \(\mathbf{pr}_1\mathbf{F}\) 。换句话说，一个对应关系 \(f = (\mathbf{F}, \mathbf{A}, \mathbf{B})\) 是一个函数，如果对于属于 f 的定义域 A 的每一个 x，关系 \((x, y) \in \mathbf{F}\) 在 y 上是函数性的（第一章，§5，第3节）；在 f 下对应于 x 的唯一对象称为 f 在 A 的元素 x 处的值，并记为 \(f(x)\) （或 \(f_x\) ，或 \(\mathbf{F}(x)\) ，或 \(\mathbf{F}_x\) ).

如果 \(f\) 是一个函数， \(F\) 是其图像， \(x\) 是 \(f\) 的定义域中的一个元素，则关系 \(y = f(x)\) 等价于 \((x, y) \in F\) (第一章，第5节，第3号，准则C46)。

备注。读者应注意，同时使用这些记号可能会引起混淆。 \(f(x)\) 并且 \(f(X)\) （与 \(f\langle X\rangle\) 同义）在定义3中引入（参见练习11）。

设 A 和 B 为两个集合；从 A 到 B 的映射是函数 f，其源（等于其定义域）等于 A，其目标等于 B；这样的函数也称为定义在 A 上且取值于 B 中的函数。

除了短语“设 f 为从 A 到 B 的映射”之外，以下短语也常被使用：“设 \(f: A \rightarrow B\) 是一个映射''，甚至“令 \(f: A \rightarrow B\) ''。为了简化涉及多个映射的论证表述，我们使用如下图表：

\[
\mathrm{A} \xrightarrow {f} \mathrm{B} \xrightarrow [ h ]{g} \mathrm{C} \xrightarrow [ j ]{i} \mathrm{E}
\]

其中，像 \(\mathbf{A} \to \mathbf{B}\) 这样的一组符号应解释为： \(f\) 是从 \(\mathbf{A}\) 到 \(\mathbf{B}\) 的映射。

若函数 \(f\) 定义在 \(\mathbf{A}\) 上，就说它把 \(x\) 变换为 \(f(x)\) （对所有 \(x \in \mathbf{A}\) ); \(f(x)\) 称为 \(x\) 在 \(f\) 下的变换，或（稍微滥用语言）称为 \(x\) 在 \(f\) 下的像。

在某些情况下，函数图被称为族；其定义域则称为指标集，而值域（通过语言的滥用）被称为族的元素集。主要正是在这种联系中，指标记号 \(f_{x}\) 用于表示f在元素x处的值。当指标集是两个集合的乘积时，我们通常称之为双重族。

同样地，目标为E的函数有时被称为E的元素族。如果E的每个元素都是集合F的子集，我们则称之为F的子集族。

在本系列中，我们经常用“函数”一词来代替“函数图”。

\minorheading{函数的例子}

\begin{enumerate}
\def\labelenumi{(\arabic{enumi})}
\item
  空集是一个函数图。每个图为空集的函数，其定义域和值域都等于空集。在这些函数中，目标集为空集的那个（即函数 \((\phi, \phi, \phi)\) ）被称为空函数。
\item
  设A为一个集合。A的恒等对应（第3号）是一个函数，称为A的恒等映射。
\end{enumerate}

因此，对于每一个集合A，都关联着一个族，它由A的恒等映射定义，其指标集为A，元素集也为A。滥用语言地，一个集合有时被称为“族”，此时它指的是与该集合相关联的那个族。

若对于 \(f\) 的定义域中的任意 \(x\) 和 \(x'\) 都有 \(f(x) = f(x')\) ，则称函数 \(f\) 为常值函数。

设 \(f\) 是集合 \(\mathbf{E}\) 到集合 \(\mathbf{E}\) 的一个映射。若 \(f(x) = x\) ，则称 \(\mathbf{E}\) 的元素 \(x\) 是 \(f\) 的不动点。

